\begin{table*}[t]
\centering
\caption{연속 위치·이산 시간 참조 모형(기간 $\kappa$ 스텝, 드론 3대, 용량 3, 행당 인스턴스 5개)과 연속 공간에서 실행한 제어기의 비교. UB·LB 는 두 차례 풀이(1시간 초기 풀이, 단축 지평·시작해로 2시간) 뒤의 incumbent 와 하한의 평균이고 ``증명''은 최적성이 증명된 인스턴스 수다. 제어기 칸은 다섯 인스턴스의 평균 makespan 과 괄호 안에 모형 UB 대비 평균 격차(\%), 행이 모두 증명되지 않았으면 LB 대비 격차도 함께 적는다(참 격차는 둘 사이). 학습 정책은 인스턴스당 30회 표본의 중앙값이며, ``전이''는 수요지 20--80·용량 5 로 학습한 정책, ``같은 분포''는 이 인스턴스와 같은 분포(수요지 8--16·용량 3)로 학습한 정책, ``증류''는 선견 규칙을 증류한 정책(제안이 아닌 분석용 변형)이다. 두 오프라인 탐색(\ref{sec:search}절)은 최적 계획을 투영의 동률 난수 10개로 다시 굴린 중앙값이다. 모든 모형 계획은 연속 환경에서 전량 배송·단절 0·추종 오차 0.0\,m 로 재생된다.}
\label{tab:cont}
\scriptsize\setlength{\tabcolsep}{1.9pt}
\begin{tabular}{rrrrrlllllll}
\toprule
 & & & & & \multicolumn{2}{c}{온라인, 학습 없음} & \multicolumn{2}{c}{오프라인 탐색} & \multicolumn{2}{c}{학습 (제안)} & 분석 \\
\cmidrule(lr){6-7}\cmidrule(lr){8-9}\cmidrule(lr){10-11}\cmidrule(lr){12-12}
$\kappa$ & $M$ & 증명 & UB & LB & 규칙 & 선견 규칙 & 방문 순서 & 작업열 & 전이 & 같은 분포 & 증류 \\
\midrule
10 & 9 & 4/5 & 130 & 118 & 199 (+54/+68) & 172 (+33/+44) & 165 (+26/+37) & 156 (+20/+31) & 235 (+83/+99) & 234 (+80/+96) & 205 (+59/+73) \\
10 & 10 & 2/5 & 154 & 118 & 269 (+78/+126) & 195 (+30/+64) & 200 (+32/+68) & 188 (+24/+59) & 268 (+75/+125) & 258 (+69/+117) & 259 (+72/+118) \\
10 & 11 & 0/5 & 166 & 122 & 261 (+55/+111) & 220 (+32/+78) & 220 (+34/+79) & 212 (+29/+73) & 279 (+67/+127) & 288 (+74/+135) & 262 (+59/+114) \\
10 & 12 & 0/5 & 210 & 118 & 285 (+34/+139) & 242 (+15/+104) & 253 (+21/+113) & 235 (+12/+98) & 306 (+45/+158) & 315 (+49/+165) & 297 (+41/+150) \\
\midrule
5 & 9 & 1/5 & 124 & 84 & 199 (+62/+134) & 172 (+39/+102) & 165 (+32/+94) & 156 (+26/+83) & 235 (+92/+176) & 234 (+89/+175) & 205 (+67/+141) \\
5 & 10 & 0/5 & 154 & 84 & 269 (+78/+218) & 195 (+30/+130) & 200 (+32/+135) & 188 (+24/+122) & 268 (+75/+215) & 258 (+69/+204) & 259 (+72/+206) \\
5 & 11 & 0/5 & 165 & 85 & 261 (+56/+203) & 220 (+32/+155) & 220 (+35/+157) & 212 (+30/+148) & 279 (+68/+225) & 288 (+75/+236) & 262 (+60/+206) \\
5 & 12 & 0/5 & 208 & 86 & 285 (+35/+228) & 242 (+16/+180) & 253 (+22/+194) & 235 (+14/+173) & 306 (+47/+254) & 315 (+51/+265) & 297 (+43/+245) \\
\bottomrule
\end{tabular}
\end{table*}
